% id: 2025-gurjao-1-1 | fonte: listas/2025/Gurjao/lista2025gurjao1.pdf, p. 1
(Physics Cup 2025) Um gás monoatómico preenche um cilindro de altura $H$ sob um pistão. O pistão oscila periodicamente para cima e para baixo com uma amplitude $a$: durante o primeiro meio período, move-se para cima com uma velocidade constante $u$, e durante o segundo meio período, move-se para baixo com a mesma velocidade constante $u$. Inicialmente, a velocidade quadrática média (RMS) das moléculas é $v$. Quanto tempo $T$ será necessário para que a velocidade RMS duplique? % sic: "monoatómico"

Utilize as seguintes hipóteses do modelo: as paredes e o pistão são isolantes térmicos perfeitos e têm capacidade térmica nula; a superfície do pistão é perfeitamente plana; o caminho livre médio das moléculas $\lambda$ satisfaz as seguintes condições: $H \gg \lambda \gg a$ e $\lambda \gg H/vt$ ; adicionalmente, $v \gg u$.
